% !TEX root = ../algebra_02.tex
\section{Ортогонализация Грама-Шмидта}

\begin{Definition}{ОНС}{}
	Система $e_1,\ldots,e_k \in V$ называется \textit{ортонормированной} (ОНС), если
	\[
	(e_i, e_j) = \delta_{ij} \quad \forall\, i, j.
	\] 
\end{Definition}

\begin{Theorem}{Ортогонализация Грама-Шмидта}{}
	Пусть $v_1, \dots, v_n$ — линейно независимая система (ЛНС) в евклидовом пространстве $V$.
	Тогда существует ортонормированная система (ОНС) $e_1, \dots, e_n$ такая, что:
	\[
	\mathrm{Lin}(e_1, \dots, e_k) = \mathrm{Lin}(v_1, \dots, v_k) \quad \forall k = 1, \dots, n.
	\]
\end{Theorem}

\begin{proof}
	Доказательство проведем по индукции.

	\textbf{База индукции} ($n = 1$):
	Положим $e_1 = \frac{v_1}{\|v_1\|}$. Очевидно, что $\|e_1\| = 1$ и $\mathrm{Lin}(e_1) = \mathrm{Lin}(v_1)$.

	\textbf{Шаг индукции} ($n = k - 1 \to n = k$):
	Пусть уже построена ОНС $e_1, \dots, e_{k-1}$ для векторов $v_1, \dots, v_{k-1}$.
	Будем искать новый вектор $e$ в виде:
	\[
	e = v_k + \alpha_1 e_1 + \dots + \alpha_{k-1} e_{k-1}
	\] 
	Потребуем, чтобы вектор $e$ был ортогонален всем предыдущим векторам базиса, то есть $(e, e_j) = 0$ для всех $j = 1, \dots, k-1$.
	Умножим скалярно обе части на $e_j$:
	\[
	(e, e_j) = (v_k, e_j) + \alpha_j (e_j, e_j) = (v_k, e_j) + \alpha_j = 0
	\] 
	Отсюда находим коэффициенты: $\alpha_j = -(v_k, e_j)$.
	Покажем, что полученный вектор $e \neq 0$. Предположим от противного, что $e = 0$.
	Тогда:
	\[
	v_k \in \mathrm{Lin}(e_1, \dots, e_{k-1}) = \mathrm{Lin}(v_1, \dots, v_{k-1})
	\] 
	Это противоречит условию линейной независимости системы $v_1, \dots, v_n$.
	Следовательно, $e \neq 0$.

	Нормируем вектор $e$, полагая $e_k = \frac{e}{\|e\|}$. Теперь $\|e_k\| = 1$, и $(e_k, e_j) = 0$ для всех $j = 1, \dots, k-1$.
	Очевидно, что $\mathrm{Lin}(e_1, \dots, e_k) = \mathrm{Lin}(v_k, e_1, \dots, e_{k-1}) = \mathrm{Lin}(v_1, \dots, v_k)$. Шаг индукции доказан.
\end{proof}
